\chapter{Materiales e Instrumental}

Para la realización de los tres experimentos descritos en este trabajo práctico se utilizó el siguiente instrumental
de laboratorio y componentes de prueba:

\section{Instrumental de Medición}
\begin{description}
    \item[Osciloscopio Digital:] RIGOL DS1052E de dos canales, ancho de banda de $50\si{\mega\hertz}$, equipado con puntas
    de prueba X10 compensadas y puntas de prueba X1.
    \item[Generador de Funciones:] Generador de señales de laboratorio con salidas senoidal, triangular y
    cuadrada con ajuste de frecuencia y amplitud.
\end{description}

\section{Componentes bajo Ensayo y de Circuito}
\begin{itemize}
    \item Bobina con núcleo magnético (hierro).
    \item Bobina con núcleo de aire montada en placa de prueba.
    \item Capacitores electrolíticos de aluminio de valores nominales: $1\uF$, $2.2\uF$ y $4.7\uF$.
    \item Resistencias serie auxiliares de limitación y sensor de corriente: $r = 100\kohm$, $r = 1\kohm$ y $r = 50\ohm$.
\end{itemize}
